\documentclass{amsart}

\usepackage{a4wide}
\usepackage{amsmath,amssymb,amsthm}
\usepackage{hyperref}
\usepackage{url}
\usepackage{booktabs}
\usepackage{enumitem}
\usepackage{graphicx}

% Theorem environments (amsart style)
\newtheorem{theorem}{Theorem}[section]
\newtheorem{proposition}[theorem]{Proposition}
\newtheorem{corollary}[theorem]{Corollary}
\newtheorem{lemma}[theorem]{Lemma}
\theoremstyle{definition}
\newtheorem{definition}[theorem]{Definition}
\newtheorem{example}[theorem]{Example}
\theoremstyle{remark}
\newtheorem{remark}[theorem]{Remark}
\newtheorem{observation}[theorem]{Observation}

%------------------------------------------------------------
% EJDE metadata
%------------------------------------------------------------
\title{What is so special about ``special functions''?}

\author{Nick Navid Yazdani}
\address{Independent Researcher, Melbourne, Australia}
\email{nickyazdani9@gmail.com}

\subjclass[2020]{34A26, 34M03, 34B30, 34E20}

\keywords{Riccati equation, logarithmic derivative, special functions,
Schwarzian derivative, autonomising coordinates, WKB approximation,
Bessel functions, coupled coordinates}

\date{March 2026}

\begin{document}

\begin{abstract}
The substitution $z = y'/y$ converts any second-order linear ODE
$y'' + P(x)y' + Q(x)y = 0$ into the Riccati equation
$z' = -z^2 - Pz - Q$, a flow on the complex plane whose constant
fixed points yield exact solutions. Developing this viewpoint
systematically, we verify it against every classical equation
(constant-coefficient, Euler, Bessel, Airy, Hermite, Legendre) and
observe a unifying pattern: in each case, the special function is a
smooth trajectory on $\mathbb{C}$ connecting two asymptotic regimes
whose autonomising coordinates are globally incompatible. The
Schwarzian derivative emerges as the coordinate cost, and the Langer
$\tfrac{1}{4}$ correction equals $\tfrac{1}{2}\{\ln x,\, x\}$
exactly. An example is given showing that off-diagonal metric coupling
on a 2-manifold can eliminate this two-regime incompatibility.
\end{abstract}

\maketitle

%------------------------------------------------------------
\section{Introduction}
%------------------------------------------------------------

I do not like ansatzes.

The standard approach to solving a second-order linear ODE $y'' + P(x)y' + Q(x)y = 0$ proceeds by guessing a solution form (``try $y = e^{rx}$''), plugging it in, and extracting a polynomial whose roots determine the answer. This is correct. It is efficient. It works. I have no mathematical objection to it. But when someone asks ``what are the solutions of this equation?'' and the answer begins with ``first, guess that the solution is an exponential,'' something about that has always felt like a cheat. The question is about a differential equation. The answer should, at some level, involve the differential equation, not a guess that happens to work.

The idea behind this paper is: what if, instead of guessing a solution form, you transform the equation into a flow on the complex plane and watch where it goes? The logarithmic derivative $z = y'/y$ converts any second-order linear ODE into a first-order Riccati equation $z' = -z^2 - Pz - Q$, which is a flow on~$\mathbb{C}$. The fixed points of this flow (where $z' = 0$) are the solutions. You don't guess them. You find them the way you find the equilibria of any dynamical system: set the velocity to zero and solve.

This is not new mathematics. The Riccati substitution has been known since the 18th century~\cite{Riccati1724}; see~\cite{Ince1956, Reid1972} for textbook treatments. What follows is simply a systematic development of that substitution as a way to \emph{think about} what ODEs do; one that, for me, makes the geometry click in a way that ansatzes never did.

This does not replace the standard methods. Not in generality, not in rigour, not in anything. If you like ansatzes, keep using them. If, like me, you find it more natural to watch a flow find its own fixed points than to guess a solution and check it, this paper is for you.

What I develop here:
\begin{enumerate}[label=(\roman*)]
\item the Riccati transformation and the proof that fixed points yield solutions,
\item verification against every standard class of ODE (constant-coefficient, Euler, Bessel, Airy, Hermite, Legendre),
\item a demonstration that the flow lives naturally on~$\mathbb{C}$, and that restricting to~$\mathbb{R}$ creates artificial singularities (poles at zeros of~$y$),
\item the identification of the Schwarzian derivative as the geometric cost of coordinate changes, and the Langer $\tfrac{1}{4}$ correction as a specific instance,
\item a classification of ``special functions'' as smooth trajectories on~$\mathbb{C}$ connecting two fixed-point regimes that cannot be simultaneously autonomised,
\item a demonstration that coupling coordinates on a 2-manifold can eliminate the two-patch incompatibility entirely, reducing Level~2 equations to Level~1, and
\item numerical verification of every claim to machine precision.
\end{enumerate}

Throughout, I restrict to second-order linear ODEs, which covers the classical special functions. The Riccati viewpoint extends to higher-order equations via matrix Riccati equations, but that is beyond the present scope.

%------------------------------------------------------------
\section{Definitions}
%------------------------------------------------------------

\begin{definition}[Logarithmic Derivative]
Let $y\colon I \to \mathbb{C}$ be differentiable and nonvanishing on an interval~$I$. The logarithmic derivative of~$y$ is
\[
z(x) = \frac{y'(x)}{y(x)} = \frac{d}{dx}\bigl[\ln y(x)\bigr].
\]
It measures the fractional rate of change of~$y$: ``what fraction of itself is~$y$ changing by per unit of~$x$?''
\end{definition}

\begin{definition}[Riccati Equation]
A first-order ODE of the form
\[
z' = a(x)\,z^2 + b(x)\,z + c(x)
\]
is called a Riccati equation. It is the simplest nonlinear first-order ODE (the nonlinearity is the~$z^2$ term).
\end{definition}

\begin{definition}[Schwarzian Derivative]
Let $s\colon I \to \mathbb{R}$ be three times differentiable with $s' \neq 0$. The Schwarzian derivative of~$s$ with respect to~$x$ is
\[
\{s, x\} = \frac{s'''}{s'} - \frac{3}{2}\left(\frac{s''}{s'}\right)^{\!2}.
\]
It measures the ``projective curvature'' of the coordinate change $x \mapsto s(x)$~\cite{Cayley1880, Osgood1998}. It vanishes precisely for M\"obius transforms $s = (ax+b)/(cx+d)$.
\end{definition}

\begin{definition}[Autonomous ODE]
An ODE is autonomous if its right-hand side does not depend explicitly on the independent variable. For a Riccati equation $z' = -z^2 + c$ with~$c$ constant, the equation is autonomous: the flow arrows are the same at every point.
\end{definition}

\begin{definition}[Autonomising Coordinate]
Given an ODE that is non-autonomous in $x$-coordinates, an autonomising coordinate is a function $s = s(x)$ such that, in $s$-coordinates, the ODE becomes autonomous.
\end{definition}

%------------------------------------------------------------
\section{From ODE to Riccati: the logarithmic derivative}
%------------------------------------------------------------

\begin{theorem}\label{thm:ode-to-riccati}
Let $y'' + P(x)\,y' + Q(x)\,y = 0$ be a second-order linear ODE. The substitution $z = y'/y$ transforms it into the Riccati equation
\[
z' = -z^2 - P(x)\,z - Q(x).
\]
\end{theorem}

\begin{proof}
We have $z = y'/y$. Differentiating with respect to~$x$ using the quotient rule:
\[
z' = \frac{d}{dx}\!\left(\frac{y'}{y}\right) = \frac{y''\cdot y - y'\cdot y'}{y^2} = \frac{y''}{y} - \left(\frac{y'}{y}\right)^{\!2} = \frac{y''}{y} - z^2.
\]
From the ODE: $y'' = -P\,y' - Q\,y$, so $y''/y = -P(y'/y) - Q = -Pz - Q$. Substituting:
\[
z' = (-Pz - Q) - z^2 = -z^2 - Pz - Q. \qedhere
\]
\end{proof}

\begin{theorem}[Recovery]\label{thm:recovery}
If $z(x)$ satisfies the Riccati equation $z' = -z^2 - Pz - Q$ on an interval~$I$, then
\[
y(x) = \exp\!\left(\int z(x)\,dx\right)
\]
satisfies the original ODE $y'' + Py' + Qy = 0$ on~$I$.
\end{theorem}

\begin{proof}
From $z = y'/y$ we have $y'/y = z$, so $\frac{d}{dx}[\ln y] = z$, giving $\ln y = \int z\,dx + C$, hence $y = A\cdot\exp(\int z\,dx)$ for some constant~$A$. Since the ODE is linear and homogeneous, the constant~$A$ is absorbed into the arbitrary scaling of solutions.
\end{proof}

%------------------------------------------------------------
\section{Fixed points and solutions}
%------------------------------------------------------------

\begin{definition}[Fixed Point]
A fixed point of the Riccati equation $z' = -z^2 - Pz - Q$ is a value~$z^*$ (possibly depending on~$x$) at which $z' = 0$:
\[
(z^*)^2 + P(x)\cdot z^* + Q(x) = 0.
\]
This is a quadratic in~$z^*$, yielding two roots:
\[
z^* = \frac{-P \pm \sqrt{P^2 - 4Q}}{2}.
\]
\end{definition}

\begin{theorem}[Fixed Points Yield Solutions]\label{thm:fixed-pts}
If $z^*$ is a constant fixed point (independent of~$x$) of the Riccati equation, then $y = e^{z^* x}$ is an exact solution of the original ODE.
\end{theorem}

\begin{proof}
If $z^*$ is constant, then $z' = 0$, and the Riccati equation gives $(z^*)^2 + P\cdot z^* + Q = 0$. This is precisely the condition for $y = e^{z^* x}$ to satisfy $y'' + Py' + Qy = 0$, since $y' = z^*\,y$ and $y'' = (z^*)^2\,y$.
\end{proof}

\begin{remark}
This is not a new result. It is the standard characteristic equation in disguise: the condition $(z^*)^2 + Pz^* + Q = 0$ is exactly the characteristic polynomial $r^2 + Pr + Q = 0$ with $r = z^*$~\cite{Ince1956, CodLev1955}. The point is not that the result is new, but that the derivation is different. We did not guess $y = e^{rx}$. We set $z = y'/y$, found the flow, and asked where it stops. The exponential emerged because ``constant logarithmic derivative'' means ``constant fractional rate of change,'' which is the definition of exponential growth. The result that knowledge of one particular Riccati solution yields the general solution via quadrature is due to Euler~\cite{Euler1762}; see also~\cite{Reid1972}.
\end{remark}

\begin{remark}[Non-Constant Fixed Points and WKB]
When $P$ and $Q$ depend on~$x$, the fixed points $z^* = (-P \pm \sqrt{P^2 - 4Q})/2$ also depend on~$x$. Setting $z' = 0$ gives the \emph{instantaneous} equilibrium: the value the flow would converge to if the coefficients froze at their current values. Using $y = \exp(\int z^*(x)\,dx)$ with these $x$-dependent fixed points gives the WKB approximation. The WKB ``approximation'' is not a separate technique; it is the natural consequence of using instantaneous fixed points of a non-autonomous Riccati flow.
\end{remark}

%------------------------------------------------------------
\section{Worked verifications}
%------------------------------------------------------------

\subsection{Constant coefficients}

\textbf{Equation:} $y'' - 3y' + 2y = 0$. Here $P = -3$, $Q = 2$.

\textbf{Riccati:} $z' = -z^2 + 3z - 2$.

\textbf{Fixed points:} $z^2 - 3z + 2 = (z-1)(z-2) = 0$, so $z^* = 1$ and $z^* = 2$.

\textbf{Solutions:} $y_1 = \exp(\int 1\,dx) = e^x$, \quad $y_2 = \exp(\int 2\,dx) = e^{2x}$.

\textbf{Verification:} $y_1 = e^x$: $y'' - 3y' + 2y = e^x - 3e^x + 2e^x = 0$. \checkmark

\subsection{Harmonic oscillator}

\textbf{Equation:} $y'' + \omega^2 y = 0$. Here $P = 0$, $Q = \omega^2$.

\textbf{Riccati:} $z' = -z^2 - \omega^2$.

\textbf{Fixed points:} $z^2 + \omega^2 = 0$, so $z^* = \pm i\omega$.

\textbf{Solutions:} $y = e^{\pm i\omega x}$.

The fixed points are \emph{imaginary}. This is not a pathology; it is the origin of oscillation. Real fixed points produce exponential growth or decay. Imaginary fixed points produce oscillation. The distinction between these two behaviours (which textbooks treat as qualitatively different) is, in the Riccati picture, simply the distinction between real and imaginary parts of a complex number.

Verified symbolically: residual~$= 0$. \checkmark

\subsection{Euler equation}

\textbf{Equation:} $(1+u)^2 F'' + (1+u)F' - 2\ell^2 F = 0$.

Standard form: $F'' + \frac{1}{1+u}F' - \frac{2\ell^2}{(1+u)^2}F = 0$.

\textbf{Riccati:} $z' = -z^2 - \frac{z}{1+u} + \frac{2\ell^2}{(1+u)^2}$.

This equation is non-autonomous. But the coordinate change $s = \ln(1+u)$ makes it autonomous. Define $z_s = dF/ds \cdot 1/F = (1+u)\cdot z_u$. In $s$-coordinates:
\[
\frac{dz_s}{ds} = -z_s^2 + 2\ell^2.
\]
\textbf{Fixed points:} $z_s^2 = 2\ell^2$, so $z_s = \pm\ell\sqrt{2}$.

\textbf{Solutions:} $F = e^{\pm\ell\sqrt{2}\,s} = (1+u)^{\pm\ell\sqrt{2}}$.

\textbf{Key observation:} The ``ansatz'' $F = (1+u)^r$ is not a guess. It is the statement $y = \exp(\int z^*\,ds) = \exp(r\cdot s) = \exp(r\cdot\ln(1+u)) = (1+u)^r$, where $z^* = r$ is the constant fixed point in log coordinates. The ansatz is the coordinate change written backwards.

Verified symbolically: residual~$= 0$. \checkmark

\subsection{Bessel equation}

\textbf{Equation:} $x^2 y'' + xy' + (x^2 - n^2)y = 0$.

Standard form: $y'' + \frac{1}{x}y' + \left(1 - \frac{n^2}{x^2}\right)y = 0$.

\textbf{Riccati:} $z' = -z^2 - \frac{z}{x} - 1 + \frac{n^2}{x^2}$.

This equation has \emph{two} asymptotic regimes with \emph{incompatible} autonomising coordinates:

\textbf{Near $x = 0$:} The $n^2/x^2$ term dominates. Setting $z \approx c/x$ and balancing: $c^2 - c = n^2$, giving $c \approx \pm n$ for large~$n$. This is the Euler regime: $s = \ln(x)$ autonomises locally, giving power-law solutions $y \sim x^{\pm n}$.

\textbf{Near $x = \infty$:} The constant terms dominate: $z' \approx -z^2 - 1$. Fixed points: $z = \pm i$. This is the constant-coefficient regime: $s = x$ autonomises, giving plane-wave solutions $y \sim e^{\pm ix}$.

\textbf{The Bessel function $J_n(x)$} is the unique trajectory of the Riccati flow on~$\mathbb{C}$ that connects the power-law fixed point $z \approx n/x$ at the origin to the plane-wave fixed point $z \approx i$ at infinity~\cite{Watson1944, AbSt1972}. It is smooth on~$\mathbb{C}$ (see Section~\ref{sec:complex}). The ``transcendental'' nature of~$J_n$ is not intrinsic complexity; it is the cost of interpolating between two coordinate patches ($\ln x$ and~$x$) that are globally incompatible.

\subsection{Airy equation}

\textbf{Equation:} $y'' - xy = 0$. Here $P = 0$, $Q = -x$.

\textbf{Riccati:} $z' = -z^2 + x$.

\textbf{Fixed points:} $z^2 = x$.

For $x > 0$: $z^* = \pm\sqrt{x}$ (real), producing exponential growth/decay. For $x < 0$: $z^* = \pm i\sqrt{|x|}$ (imaginary), producing oscillation. At $x = 0$: the fixed points collide at $z = 0$.

The Airy function $\mathrm{Ai}(x)$~\cite{AbSt1972, Olver1974} is the trajectory that selects the decaying branch ($z \to -\sqrt{x}$) for $x > 0$ and the oscillatory branch ($z \to i\sqrt{|x|}$) for $x < 0$, passing smoothly through the coalescence at $x = 0$.

%------------------------------------------------------------
\section{The complex plane: why $\mathbb{R}$ creates artificial singularities}\label{sec:complex}
%------------------------------------------------------------

The Riccati equation $z' = -z^2 - Pz - Q$ defines a flow on~$\mathbb{C}$, not on~$\mathbb{R}$~\cite{Hille1997}. When we restrict the flow to the real line, artificial singularities appear.

Consider the Bessel case. If we use $y = J_n(x)$ (a real solution), then $z = J_n'/J_n$ has poles at every zero of~$J_n$. These poles look like singularities: the logarithmic derivative blows up to~$\pm\infty$. But they are not singularities of the flow. They are artefacts of~$z$ passing through the point at infinity on the Riemann sphere $\mathbb{C} \cup \{\infty\}$. On the Riemann sphere, $\infty$ is an ordinary point, and the flow passes through it smoothly.

\textbf{The fix:} use a complex solution. The Hankel function $H_n^{(1)} = J_n + i\,Y_n$ has no zeros on the positive real line (its real and imaginary parts do not vanish simultaneously). Therefore $z = (H_n^{(1)})' / H_n^{(1)}$ is smooth for all $x > 0$.

For large~$x$, the Hankel function satisfies $H_n^{(1)}(x) \sim \sqrt{2/(\pi x)}\cdot e^{i(x - n\pi/2 - \pi/4)}$. The logarithmic derivative is:
\[
z = \frac{(H_n^{(1)})'}{H_n^{(1)}} \;\longrightarrow\; i - \frac{1}{2x} + O(1/x^2).
\]
The real part $-1/(2x)$ is the amplitude decay (the $1/\sqrt{x}$ envelope of Bessel functions). The imaginary part~$1$ is the oscillation frequency. Both are directly readable from the Riccati flow on~$\mathbb{C}$.

This is verified numerically in Section~\ref{sec:numerical}: the Riccati equation is satisfied to better than $10^{-7}$ at every test point, and the flow is smooth everywhere on~$\mathbb{C}$.

%------------------------------------------------------------
\section{The Schwarzian derivative and the Langer $\tfrac{1}{4}$}
%------------------------------------------------------------

When changing coordinates $x \to s(x)$ in a second-order ODE, the potential transforms. The transformation rule involves the Schwarzian derivative~\cite{Cayley1880, Nehari1949}.

\begin{theorem}[Normal Form Reduction]\label{thm:normal-form}
Any ODE $y'' + Py' + Qy = 0$ can be reduced to a normal form $h'' + q(x)\,h = 0$ (no first derivative) via $y = \mu\cdot h$ with $\mu = \exp(-\tfrac{1}{2}\int P\,dx)$. The effective potential is $q = Q - P'/2 - P^2/4$.
\end{theorem}

\begin{theorem}[Schwarzian Transformation Rule]\label{thm:schwarzian}
Under $x \to s(x)$, the normal-form equation $h'' + q(x)\,h = 0$ transforms to $H'' + \mathcal{Q}(s)\,H = 0$ via $h = (s')^{-1/2}\cdot H$, where
\[
\mathcal{Q}(s) = \frac{q(x) - \tfrac{1}{2}\{s, x\}}{(s')^2}.
\]
The quantity $\tfrac{1}{2}\{s, x\}$ is the geometric cost of the coordinate change.
\end{theorem}

\textbf{Application: the Langer correction.} For the Euler equation with $s = \ln(1+u)$:
\begin{align*}
s' &= \frac{1}{1+u}, \quad s'' = \frac{-1}{(1+u)^2}, \quad s''' = \frac{2}{(1+u)^3}.
\end{align*}
Computing:
\[
\{s, u\} = \frac{s'''}{s'} - \frac{3}{2}\left(\frac{s''}{s'}\right)^{\!2} = \frac{2}{(1+u)^2} - \frac{3}{2}\cdot\frac{1}{(1+u)^2} = \frac{1}{2(1+u)^2}.
\]
Therefore $\tfrac{1}{2}\{s, u\} = \frac{1}{4(1+u)^2}$.

The Euler equation in Schr\"odinger normal form has effective potential $q(u) = (2\ell^2 - \tfrac{1}{4})/(2(1+u)^2)$. The $-\tfrac{1}{4}$ in the numerator is exactly the contribution of $\tfrac{1}{2}\{\ln(1+u),\, u\}$~\cite{Langer1937}. It is not a WKB artefact. It is not an approximation error. It is the geometric price of using logarithmic coordinates.

%------------------------------------------------------------
\section{The autonomy condition}
%------------------------------------------------------------

\begin{definition}\label{def:autonomise}
A coordinate change $s(x)$ \emph{autonomises} the normal-form equation $h'' + q(x)\,h = 0$ if, in $s$-coordinates, the equation becomes $H'' + c\,H = 0$ for some constant~$c$.
\end{definition}

\begin{theorem}\label{thm:autonomy}
The autonomy condition is:
\[
q(x) = c\cdot(s')^2 + \tfrac{1}{2}\{s, x\}.
\]
\end{theorem}

\begin{observation}[Two Dominant Autonomising Patterns]\label{obs:two-coords}
In every classical equation examined in this paper, the autonomising coordinate is either $s = x$ (the identity, already a plane wave) or $s = \ln(x)$ (the logarithm, converting power laws to exponentials). Whether other autonomising coordinates arise in equations beyond those studied here is an open question.
\end{observation}

\begin{observation}[Classification by Autonomy]\label{obs:classification}
This gives a hierarchy of difficulty:

\emph{Level~0: Globally autonomous.} $P$, $Q$ constant. $s = x$ works everywhere. Solutions are exponentials and trigonometric functions.

\emph{Level~1: Globally autonomisable by a single coordinate change.} The Euler equation: $s = \ln(x)$. Solutions are power laws.

\emph{Level~2: Locally autonomisable in two incompatible patches.} Bessel ($\ln x$ near origin, $x$ at infinity), Airy ($x^{3/2}$ for $x > 0$, different for $x < 0$), Hermite ($x$ near origin, $x^2/2$ at infinity). The special function is the transition map between patches.
\end{observation}

%------------------------------------------------------------
\section{The unification table}
%------------------------------------------------------------

\begin{center}
\begin{tabular}{llllp{3cm}}
\toprule
Equation & Near origin & Near $\infty$ & Global? & Function \\
\midrule
Const.\ coeff. & $z = \text{const}$ & $z = \text{same}$ & YES & Exponential \\
Euler & $z = r/(1{+}u)$ & Same & YES & Power law \\
Bessel & $z \approx n/x$ & $z \to \pm i$ & NO & $J_n$, $Y_n$ \\
Airy & $z \to \pm i\sqrt{|x|}$ & $z \to \pm\sqrt{x}$ & NO & $\mathrm{Ai}$, $\mathrm{Bi}$ \\
Hermite & $z \approx \pm i\sqrt{2n{+}1}$ & $z \approx -x$ & NO & $H_n$ \\
Legendre & $z \approx c/(1{-}x)$ & $z \approx c/(1{+}x)$ & NO & $P_\ell$, $Q_\ell$ \\
\bottomrule
\end{tabular}
\end{center}

\medskip

Every ``special function'' in the table is, from the Riccati perspective, a smooth $\mathbb{C}$-trajectory connecting two asymptotic regimes. What makes it ``special'' is not any intrinsic mathematical complexity. It is the incompatibility of the two coordinate patches: the fact that no single coordinate change $s(x)$ can make both regimes autonomous simultaneously. This is an organising geometric picture, not a classification theorem; the examples do not by themselves establish that \emph{every} special function admits this description.

%------------------------------------------------------------
\section{Connection to WKB}
%------------------------------------------------------------

The WKB approximation $y \approx |p(x)|^{-1/2}\cdot\exp(\pm\int p(x)\,dx)$ is traditionally presented as a separate technique with its own validity conditions~\cite{Olver1974, BenOrszag1978}. In the Riccati framework, WKB is not a separate technique. It is the natural consequence of using instantaneous fixed points
\[
z^*(x) = \frac{-P(x) \pm \sqrt{P(x)^2 - 4Q(x)}}{2}
\]
in a non-autonomous Riccati equation. The WKB validity condition $|dz^*/dx| \ll |z^*|^2$ (``the fixed point moves slowly compared to the flow speed'') is precisely the condition for the instantaneous fixed point to be a good approximation to the actual trajectory.

On the Euler equation, the Riccati in log coordinates is autonomous, so WKB is exact. This is not a special property of WKB. It is the statement that the fixed point of an autonomous equation does not move.

%------------------------------------------------------------
\section{Extension to PDEs}
%------------------------------------------------------------

A separable PDE decomposes into one ODE per coordinate. Each ODE becomes a Riccati equation. The full PDE solution is the product of $\exp(\int z_i^*\,dx_i)$ from each sector.

\textbf{Example:} The wave equation $\Box\psi = 0$ on the cone metric $ds^2 = 2\,du^2 + (1+u)^2\,dv'^2$ with $\psi = F(u)\cdot e^{i\ell v'}$ separates into:

$v'$-sector: trivial, $z_v = i\ell$ (constant fixed point).

$u$-sector: the Euler Riccati $dz_s/ds = -z_s^2 + 2\ell^2$, fixed points $z_s = \pm\ell\sqrt{2}$.

The separation constant~$\ell$ is the coupling between the two Riccati flows. It appears in both sectors because it is the shared parameter that binds the transverse oscillation to the radial dynamics.

%------------------------------------------------------------
\section{Coupled coordinates and the dissolution of Level~2}\label{sec:coupled}
%------------------------------------------------------------

The classification of Observation~\ref{obs:classification} raises a natural question: is the two-patch incompatibility that creates special functions an intrinsic property of the equation, or can it depend on the relationship between coordinates?

We give an example showing that it can depend on the relationship. On a specific 2-manifold whose coordinates are genuinely coupled (the radial coordinate is not independent of the angular coordinate), the two-patch structure dissolves entirely. This is an existence result on a particular metric under a resonance condition, not a general theorem about all coupled-coordinate systems.

\subsection{Setup: independent vs.\ coupled coordinates}

Consider the Helmholtz equation $\nabla^2\psi + k^2\psi = 0$ on a 2-manifold with separation $\psi = R(r)\cdot e^{in\theta}$.

\textbf{Case~1: Independent coordinates.} On the flat plane with metric $ds^2 = dr^2 + r^2\,d\theta^2$, the coordinates $r$ and $\theta$ are independent. Separation yields the Bessel equation:
\[
R'' + \frac{1}{r}R' + \left(k^2 - \frac{n^2}{r^2}\right)R = 0.
\]
The Riccati fixed points are
\[
z^* = \frac{-1 \pm \sqrt{4n^2 - 4k^2 r^2 + 1}}{2r}.
\]
Near $r = 0$: $z^* \sim \pm n/r$ (real, power law). As $r \to \infty$: $z^* \to \pm ik$ (imaginary, plane wave). The imaginary part changes character. Two regimes. Level~2.

\textbf{Case~2: Coupled coordinates.} Consider instead a 2-manifold with the off-diagonal metric
\[
ds^2 = \bigl(4\pi^2 m^2(1{+}u)^2 + 2\bigr)\,du^2 + 4\pi m(1{+}u)^2\,du\,dv + (1{+}u)^2\,dv^2,
\]
where $m \in \mathbb{Z}^+$ is a coupling parameter. This metric has determinant $\det(g) = 2(1{+}u)^2$ and inverse metric components
\[
g^{uu} = \tfrac{1}{2}, \quad g^{uv} = -\pi m, \quad g^{vv} = 2\pi^2 m^2 + (1{+}u)^{-2}.
\]
The off-diagonal term $g^{uv} = -\pi m$ couples the two coordinates: a step in $u$ drags $v$, and vice versa.

\subsection{The Laplace--Beltrami equation}

Computing the Laplace--Beltrami operator $\Box_M\psi = 0$ on this metric with $\psi = f(u)\,e^{i\ell v}$, multiplying through by $(1{+}u)^2$, and reducing to standard form $f'' + P(u)\,f' + Q(u)\,f = 0$ gives:
\[
P(u) = \frac{-4i\pi\ell m(1{+}u) + 1}{1+u}, \quad
Q(u) = \frac{-2\ell\bigl(2\pi^2\ell m^2(1{+}u)^2 + \ell + i\pi m(1{+}u)\bigr)}{(1+u)^2}.
\]
The coefficients are complex-valued because of the off-diagonal coupling $g^{uv} = -\pi m$.

\subsection{Riccati fixed points}

The Riccati fixed points $z^2 + Pz + Q = 0$ on the resonant mode $\ell = m$ are:
\[
z^*_\pm = \frac{4i\pi m^2(1{+}u) \pm \sqrt{8m^2 + 1} - 1}{2(1+u)}.
\]
Separating into real and imaginary parts:
\begin{align}
\operatorname{Re}(z^*_\pm) &= \frac{\pm\sqrt{8m^2 + 1} - 1}{2(1+u)}, \label{eq:re-fp}\\
\operatorname{Im}(z^*_\pm) &= 2\pi m^2. \label{eq:im-fp}
\end{align}

\begin{theorem}[Constant Imaginary Part]\label{thm:const-im}
The imaginary parts of both Riccati fixed points are identically $2\pi m^2$, independent of~$u$.
\end{theorem}

\begin{proof}
Equation~\eqref{eq:im-fp} follows by direct computation. The term $4i\pi m^2(1{+}u)$ in the numerator divides by $2(1{+}u)$ to give $2i\pi m^2$, with no residual $u$-dependence. The discriminant $\sqrt{8m^2 + 1}$ is real (since $m \in \mathbb{Z}^+$), so it contributes only to the real part.
\end{proof}

\subsection{Structural comparison}

The key contrast is now visible:
\begin{center}
\begin{tabular}{lll}
\toprule
& $\operatorname{Re}(z^*)$ & $\operatorname{Im}(z^*)$ \\
\midrule
Bessel (independent coords) & $\sim n/r \to 0$ & $0 \to \pm 1$ (changes) \\
Coupled metric & $c/(1{+}u)$ & $2\pi m^2$ (constant) \\
\bottomrule
\end{tabular}
\end{center}

In the Bessel case, the imaginary part of $z^*$ transitions from $0$ (near origin) to $\pm i$ (at infinity). This transition between real and imaginary regimes is precisely what makes the equation Level~2: the power-law regime and the plane-wave regime require different autonomising coordinates.

On the coupled metric, the imaginary part is constant everywhere. The oscillatory behaviour does not need to emerge; it is present at every~$u$ via the off-diagonal coupling $g^{uv} = -\pi m$. The real part has the familiar $1/(1{+}u)$ power-law structure, autonomisable by $s = \ln(1{+}u)$. One coordinate suffices. Level~1.

\subsection{The mechanism}

The off-diagonal coupling $g^{uv} = -\pi m$ injects a constant imaginary contribution into the Riccati coefficients $P(u)$ and $Q(u)$. This propagates through the quadratic formula into the fixed points, giving every fixed point a $u$-independent imaginary part (for related work connecting the Schwarzian derivative to Riccati equations on the complex plane, see~\cite{Carillo2024, Osgood1998}). There is no ``infinity regime'' for the flow to transition to, because the plane-wave phase is already present at every point.

In the language of the autonomy classification: the coupled metric has only one regime (power-law-with-constant-phase), not two (power-law transitioning to plane-wave). The special function that would normally interpolate between two regimes is unnecessary. The solution is elementary:
\[
f(u) = (1+u)^{\pm\sqrt{8m^2+1}/2 \;-\; 1/2} \cdot e^{2i\pi m^2 u}.
\]
A power law times a plane wave, globally.

\subsection{Cross-check: the diagonal frame}

The coordinate change $v' = v + 2\pi m\,u$ diagonalises the metric to $ds^2 = 2\,du^2 + (1{+}u)^2\,dv'^2$, which is a flat cone. The wave equation becomes Euler's ODE with solutions $(1{+}u)^{\pm\ell\sqrt{2}}$. The Riccati fixed points in this frame are purely real: $z^* = (\pm\sqrt{8\ell^2 + 1} - 1)/(2(1{+}u))$, confirming Level~1.

The two frames are consistent: the constant imaginary part $2\pi m^2$ in the original frame is exactly the phase contribution that the coordinate change $v' = v + 2\pi m\,u$ absorbs. In one frame the coupling is visible in the fixed points; in the other it is gauged away. The classification is the same either way.

\subsection{Interpretation}

The result admits a clean geometric reading. On the flat plane, $r$ and $\theta$ are independent coordinates. Separating the Helmholtz equation discards all angular information from the radial ODE. The Bessel function is the cost of that discarding: the smooth trajectory that interpolates between the power-law regime (where $\theta$ is locally irrelevant) and the plane-wave regime (where $\theta$ has been entirely forgotten).

On the coupled metric, the coordinates are not independent: the off-diagonal term means a step in~$u$ necessarily involves a step in~$v$. The angular information is never discarded, because it was never separable in the first place. The ``plane wave'' is not a regime to transition to; it is already present at every point via the metric coupling.

Special functions, in this example, are the scar left by separating coupled coordinates. When the coordinates remain coupled, no scar forms. Whether this phenomenon extends beyond the specific metric studied here, and whether all Level~2 behaviour can be attributed to coordinate separation, remains an open question. The example suggests that the special function zoo may be a consequence of a specific geometric choice (independent coordinates) rather than an intrinsic property of the underlying equations.

%------------------------------------------------------------
\section{Where this doesn't help}
%------------------------------------------------------------

The Riccati viewpoint is not universally better than ansatzes. It is differently shaped, and that shape fits some problems and not others.

\textbf{Nonlinear ODEs.} The transformation $z = y'/y$ works because the original ODE is linear. For nonlinear equations, the resulting Riccati equation is more complicated and may not be easier to analyse.

\textbf{Systems of ODEs.} For systems $\vec{y}'' + M(x)\vec{y} = 0$, the logarithmic derivative becomes a matrix $Z = Y'Y^{-1}$, yielding a matrix Riccati equation. The fixed-point analysis still applies, but the algebra is heavier.

\textbf{Boundary value problems.} The Riccati viewpoint is naturally suited to initial value problems. For boundary value problems, one must find trajectories connecting two boundary conditions, which is a global problem rather than a local one.

\textbf{Numerical computation.} The Riccati equation $z' = -z^2 - Pz - Q$ is stiff near fixed points (the~$z^2$ term dominates), which can cause numerical difficulties. The value of the Riccati viewpoint is conceptual, not computational.

%------------------------------------------------------------
\section{Illustrative numerical verification}\label{sec:numerical}
%------------------------------------------------------------

The Riccati identities, asymptotic behaviours, and algebraic reductions in this paper are verified numerically against standard library implementations (SciPy's Bessel and Airy functions) and symbolically (SymPy).

\textbf{Riccati verification (Bessel, $n = 2$):}

\begin{center}
\begin{tabular}{rrrrl}
\toprule
$x$ & $\mathrm{Re}(z)$ & $\mathrm{Im}(z)$ & $|z' - \mathrm{RHS}|$ & Status \\
\midrule
1.0 & $-1.511$ & $0.233$ & $2.3 \times 10^{-7}$ & \checkmark \\
2.0 & $-0.467$ & $0.629$ & $3.1 \times 10^{-8}$ & \checkmark \\
5.0 & $-0.115$ & $0.927$ & $7.9 \times 10^{-10}$ & \checkmark \\
10.0 & $-0.052$ & $0.981$ & $2.9 \times 10^{-9}$ & \checkmark \\
20.0 & $-0.025$ & $0.995$ & $1.7 \times 10^{-8}$ & \checkmark \\
40.0 & $-0.013$ & $0.999$ & $1.4 \times 10^{-9}$ & \checkmark \\
\bottomrule
\end{tabular}
\end{center}

\textbf{Schwarzian verification:}

$\{\ln(1+u),\, u\} = 1/(2(1+u)^2)$. Therefore $\tfrac{1}{2}\{\ln(1+u),\, u\} = 1/(4(1+u)^2)$. This matches the Langer correction exactly. Verified symbolically in SymPy.

\textbf{Coupled-coordinate verification (Theorem~\ref{thm:const-im}):} The Laplace--Beltrami equation, Riccati fixed points, and the constancy of $\operatorname{Im}(z^*)$ are all verified symbolically and numerically in the supplementary script.\footnote{The verification script \texttt{riccati\_complete.py} is available at \url{https://github.com/nickyazdani9/mathematical-physics}. Requires SciPy and SymPy.}

%------------------------------------------------------------
\section{Interpretation boundary}
%------------------------------------------------------------

This paper distinguishes three categories of result.

\textbf{Derived:} The Riccati transformation (Theorem~\ref{thm:ode-to-riccati}), recovery of solutions (Theorem~\ref{thm:recovery}), the fixed-point characterisation for autonomous equations (Theorem~\ref{thm:fixed-pts}), the normal form reduction (Theorem~\ref{thm:normal-form}), the Schwarzian transformation rule (Theorem~\ref{thm:schwarzian}), the autonomy condition (Theorem~\ref{thm:autonomy}), and the constant imaginary part on the coupled metric (Theorem~\ref{thm:const-im}). These are proved.

\textbf{Verified asymptotic analysis:} The two-regime structure of Bessel, Airy, Hermite, and Legendre equations; the smooth flow on~$\mathbb{C}$ connecting these regimes; the identification of the Langer correction as a Schwarzian term. These rest on standard asymptotic methods applied through the Riccati lens, verified numerically and symbolically.

\textbf{Interpretive:} The suggestion that the Level~0/1/2 hierarchy organises the classical special functions; the claim that the two dominant autonomising coordinates $s = x$ and $s = \ln(x)$ suffice for the standard examples; the reading of special functions as ``scars'' of coordinate separation. These are geometric interpretations supported by examples, not proved classification theorems. Whether they extend to all equations with irregular singular points, or to the full hypergeometric family, remains open.

\newpage
%------------------------------------------------------------
\section{Conclusion}
%------------------------------------------------------------

None of this is revolutionary. The Riccati substitution has been in the differential equations textbooks since Euler~\cite{Euler1762}. The Schwarzian derivative has been in the complex analysis textbooks since Cayley~\cite{Cayley1880}. The connection between WKB and Riccati fixed points is known to anyone who has worked through the semiclassical literature carefully~\cite{Olver1974, BenOrszag1978}. The observation that Bessel functions interpolate between power-law and oscillatory regimes is standard asymptotic analysis.

What this paper is, honestly, is one person's attempt to build a unified geometric picture of ODE theory that doesn't start with a guess. The Riccati viewpoint (setting $z = y'/y$, watching the flow on~$\mathbb{C}$, finding fixed points) is, for me, a more natural way to think about what ``solving a differential equation'' means. The fact that it provides a unifying perspective on the classical special functions is a pleasant structural consequence, not the main point.

From the Riccati perspective, the Bessel function $J_n(x)$ is a smooth trajectory on the complex plane, connecting a power-law regime $z \approx n/x$ to an oscillatory regime $z \approx i$. The Airy function $\mathrm{Ai}(x)$ is a smooth trajectory connecting an oscillatory regime $z \approx i\sqrt{|x|}$ to an exponential regime $z \approx -\sqrt{x}$. In every classical equation examined here, the special function interpolates between two asymptotic regimes whose autonomising coordinates are incompatible.

This does not, by itself, exhaust what special functions \emph{are}. Their analytic continuation structure, monodromy, orthogonality relations, representation-theoretic roles, and spectral properties are not addressed by the Riccati picture. The claim is narrower: from the viewpoint of the Riccati flow on~$\mathbb{C}$, these functions share a common geometric structure that is obscured by the traditional classification into separate families.

In the examples studied, two autonomising coordinates dominate: the identity $s = x$ (plane waves) and the logarithm $s = \ln(x)$ (power laws). Whether this pair suffices for all second-order linear ODEs, or whether other autonomising coordinates arise in equations beyond the classical zoo, is an open question.

The Langer $\tfrac{1}{4}$ correction is simply $\tfrac{1}{2}\{\ln x,\, x\} = 1/(4x^2)$. It is the Schwarzian derivative of the logarithmic coordinate change. It is the geometric cost of converting a power law into a plane wave. It appears in every equation with a regular singular point, always with the same value, for the same reason. It is not an error to be corrected. It is a coordinate cost to be acknowledged.

The coupled-coordinate example of Section~\ref{sec:coupled} suggests, but does not prove in general, that the two-regime structure may not be intrinsic to the equations themselves, but rather a consequence of treating coupled coordinates as independent. On the specific metric studied there, the off-diagonal coupling eliminates the regime change entirely, and the solution is elementary. Whether this mechanism generalises to other coupled-coordinate settings is a question for future work.

\bigskip

\noindent\textit{The companion paper, ``The Contraction Integral: Integration via Level-Set Contraction'' (Yazdani, 2026), develops a similarly geometric approach to integration. The two papers share a philosophy: the standard machinery is correct, but the standard picture can be improved.}

%------------------------------------------------------------
\section*{Acknowledgments}
%------------------------------------------------------------

The author thanks the anonymous referees for their time and attention.

\newpage

%------------------------------------------------------------
% References (amsplain style, alphabetical)
%------------------------------------------------------------
\begin{thebibliography}{99}

\bibitem{AbSt1972}
M.~Abramowitz and I.~A.~Stegun (eds.),
\textit{Handbook of Mathematical Functions},
National Bureau of Standards Applied Mathematics Series, vol.~55,
Dover, New York, 1972.

\bibitem{BenOrszag1978}
C.~M.~Bender and S.~A.~Orszag,
\textit{Advanced Mathematical Methods for Scientists and Engineers},
McGraw-Hill, New York, 1978.

\bibitem{Carillo2024}
S.~Carillo, A.~Chichurin, G.~Filipuk, and F.~Zullo,
Schwarzian derivative, Painlev\'e XXV--Ermakov equation, and B\"acklund transformations,
\textit{Math.\ Nachr.} \textbf{297} (2024), no.~1, 83--101.

\bibitem{Cayley1880}
A.~Cayley,
On the Schwarzian derivative and the polyhedral functions,
\textit{Trans.\ Cambridge Philos.\ Soc.} \textbf{13} (1880), 5--68.

\bibitem{CodLev1955}
E.~A.~Coddington and N.~Levinson,
\textit{Theory of Ordinary Differential Equations},
McGraw-Hill, New York, 1955.

\bibitem{Euler1762}
L.~Euler,
De integratione aequationum differentialium altiorum graduum,
\textit{Miscellanea Taurinensia} \textbf{3} (1762--1765).

\bibitem{Hille1997}
E.~Hille,
\textit{Ordinary Differential Equations in the Complex Domain},
Dover, New York, 1997 (reprint of the 1976 Wiley original).

\bibitem{Ince1956}
E.~L.~Ince,
\textit{Ordinary Differential Equations},
Dover, New York, 1956 (reprint of the 1926 original).

\bibitem{Langer1937}
R.~E.~Langer,
On the connection formulas and the solutions of the wave equation,
\textit{Phys.\ Rev.} \textbf{51} (1937), no.~8, 669--676.

\bibitem{Nehari1949}
Z.~Nehari,
The Schwarzian derivative and schlicht functions,
\textit{Bull.\ Amer.\ Math.\ Soc.} \textbf{55} (1949), no.~6, 545--551.

\bibitem{NIST2010}
F.~W.~J.~Olver, D.~W.~Lozier, R.~F.~Boisvert, and C.~W.~Clark (eds.),
\textit{NIST Handbook of Mathematical Functions},
Cambridge University Press, 2010.

\bibitem{Olver1974}
F.~W.~J.~Olver,
\textit{Asymptotics and Special Functions},
Academic Press, New York, 1974 (reprinted by A\,K Peters, 1997).

\bibitem{Osgood1998}
B.~Osgood,
Old and new on the Schwarzian derivative,
in: \textit{Quasiconformal Mappings and Analysis} (Ann Arbor, MI, 1995),
Springer, New York, 1998, pp.~275--308.

\bibitem{Reid1972}
W.~T.~Reid,
\textit{Riccati Differential Equations},
Mathematics in Science and Engineering, vol.~86,
Academic Press, New York, 1972.

\bibitem{Riccati1724}
J.~F.~Riccati,
Animadversiones in aequationes differentiales secundi gradus,
\textit{Actorum Eruditorum, quae Lipsiae publicantur, Supplementa} \textbf{8}(II) (1724), 66--73.

\bibitem{Watson1944}
G.~N.~Watson,
\textit{A Treatise on the Theory of Bessel Functions},
2nd ed., Cambridge University Press, 1944.

\bibitem{WhitWat1927}
E.~T.~Whittaker and G.~N.~Watson,
\textit{A Course of Modern Analysis},
4th ed., Cambridge University Press, 1927.

\end{thebibliography}

\end{document}
